\documentclass[10pt,a4paper]{amsart}
\usepackage[margin=19mm]{geometry}
\usepackage{amsmath,amssymb,mathtools}
\usepackage[scr=rsfs,cal=euler]{mathalfa}
\usepackage{xfrac}
\usepackage[unicode,hidelinks,bookmarks=false]{hyperref}
\newcommand{\ie}{i.e.\ }
\newcommand{\cf}{cf.\ }
\newcommand{\resp}{respectively\ }
\newcommand{\Subsection}{\subsection}
\providecommand{\nobreakdash}{\nobreak\hbox{-}}
\setlength{\emergencystretch}{2em}
\multlinegap0pt
\usepackage{bm}
\usepackage{bbm}
\usepackage{dsfont}
\usepackage{xspace}
\usepackage{cancel}
\usepackage{mathtools}
\usepackage{cleveref}
\usepackage{amsthm}
\makeatletter
\@ifundefined{lemma}{\newtheorem{lemma}{Lemma}}{}
\@ifundefined{proposition}{\newtheorem{proposition}[lemma]{Proposition}}{}
\makeatother
\newcommand\D{\Delta}
\title[On set-theoretic solutions of pentagon equation and positive]{On set-theoretic solutions of pentagon equation and positive basis Hopf algebras}
\author{Ilaria Colazzo, Geoffrey Janssens}
\date{Source: arXiv:2601.22089v1 (CC BY 4.0)}
\begin{document}
\maketitle

\noindent\emph{Source and licence.} On set-theoretic solutions of pentagon equation and positive basis Hopf algebras. Version arXiv:2601.22089v1 (CC BY 4.0), distributed under the Creative Commons Attribution 4.0 International licence, as recorded in the arXiv licence metadata for this exact version and shown on the abstract page https://arxiv.org/abs/2601.22089v1. Full rights notice: licenses/arxiv-2601.22089-CC-BY-4.0.txt.\par\smallskip

\noindent\textit{Editorial scope.} Counted unit: the Proposition whose source label is \texttt{prop:ThetaInv} in the article (source0.tex lines 1762--1768, source label \texttt{prop:ThetaInv}), delivered together with the article's own complete original proof of it (source0.tex lines 1804--1903, one \texttt{proof} environment of 100 lines). No mathematics is written, repaired or re-derived: statement and proof are the article's own text line for line. The proof is detached from the statement in the source: the article states the result at lines 1762--1768 and prints its proof at lines 1804--1903, and the proof environment's own header names the statement, so the association is the article's own. Required context retained uncounted: eq:ThetaDef (lines 1756--1758). Every definition, display and label of those lines is retained unchanged. Omitted from this package and not redistributed: 1--1755, 1759--1761, 1769--1803, 1904--2565. Nothing inside a retained excerpt is omitted or altered: every retained body is delivered exactly as the article's lines are written, so no omission receipt of any kind accompanies this package. Citations: the retained excerpts carry no citation command at all, so no bibliography entry of the article is transcribed and no reference is redistributed. Reference adaptation: none was needed and none was made; every reference inside the delivered unit resolves inside it, so no unresolved reference or citation is produced. Printed slips kept literal: no printed word, symbol, hypothesis or number of the counted statement or of its proof was altered. Layout and class: the article's own class is \texttt{amsart} with packages including inputenc, lmodern, fontenc, microtype, xy, amsmath, amssymb, amsfonts, amsthm, latexsym, mathrsfs, mathtools, stmaryrd, xcolor; the wrapper is the lane's pinned \texttt{amsart} editorial wrapper at 10pt on A4, which restates the theorem-like environments the retained text needs and adds the article's own macro definitions that the retained text needs (\texttt{\textbackslash D}), with extra wrapper packages (bm, bbm, dsfont, xspace, cancel, mathtools, cleveref). The delivered numbering is the unit's own. Assets: no retained span contains a figure environment, a graphics-inclusion command, a PDF-page inclusion, a table or a TikZ picture; no image file is copied into this package, the package declares no asset dependency and no image file is redistributed with it. Licence: version arXiv:2601.22089v1 of this article is distributed under the Creative Commons Attribution 4.0 International licence, as recorded in the arXiv licence metadata for this exact version and shown on the abstract page https://arxiv.org/abs/2601.22089v1. A full rights notice is delivered at licenses/arxiv-2601.22089-CC-BY-4.0.txt. Characters: every retained excerpt is pure ASCII, so the delivered engine, XeLaTeX with its default Latin Modern text font, reports no missing character.
\begin{equation}\label{eq:ThetaDef}
\Theta_A:k[A^\vee]^*\longrightarrow k[A]: \delta_\chi \mapsto \frac{1}{|A|}\sum_{a\in A}\chi(a^{-1})\,a.
\end{equation}

\begin{proposition}\label{prop:ThetaInv}
The map $\Theta_A$ is an isomorphism of Hopf algebras. Its inverse $\Theta_A^{-1}:k[A]\to k[A^\vee]^*$ is given on group elements $a \in A$ by
\begin{equation}\label{eq:ThetaInv}
\Theta_A^{-1}(a)=\sum_{\chi\in A^\vee}\chi(a)\,\delta_\chi,
\end{equation}
extended $k$-linearly.
\end{proposition}

\begin{proof}[Proof of \Cref{prop:ThetaInv}]
We will regularly use the character orthogonality relations which hold due to our assumptions on the ground field $k$:
\begin{equation}\label{eq:orth}
\sum_{\chi\in A^\vee}\chi(x)=
\begin{cases}
|A|,& x=1,\\
0,& x\neq 1,
\end{cases}
\qquad
\sum_{a\in A}\chi(a)=
\begin{cases}
|A|,& \chi=1,\\
0,& \chi\neq 1.
\end{cases}
\end{equation}

\medskip\noindent
{\it Algebra map.} By definition 
$$ \Theta_A(\delta_\chi)\Theta_A(\delta_\psi)
=\frac1{|A|^2}\sum_{a,b\in A}\chi(a^{-1})\psi(b^{-1})\,ab.
$$
The coefficient of a fixed $t\in A$ equals
\[
\frac1{|A|^2}\sum_{a\in A}\chi(a^{-1})\psi\bigl((a^{-1}t)^{-1}\bigr)
=\frac{\psi(t^{-1})}{|A|^2}\sum_{a\in A}(\chi^{-1}\psi)(a).
\]
By \eqref{eq:orth} this is $\frac1{|A|}\chi(t^{-1})$ if $\chi=\psi$ and $0$ otherwise. Hence
$\Theta_A(\delta_\chi)\Theta_A(\delta_\psi)=\delta_{\chi,\psi}\Theta_A(\delta_\chi)=\Theta_A(\delta_\chi\delta_\psi)$.
The orthogonality relations also yield
\[
\Theta_A(1)=\Theta_A\Bigl(\sum_{\chi}\delta_\chi\Bigr)
=\frac1{|A|}\sum_{a\in A}\Bigl(\sum_{\chi}\chi(a^{-1})\Bigr)a
=1.
\]


\medskip\noindent
{\it Coalgebra map.}
We show $\D_{k[A]} (\Theta_A(\delta_\chi))=(\Theta_A\otimes \Theta_A)\D_{k[A^{\vee}]^*} (\delta_\chi)$ and drop the indices of the coproducts.
The left-hand side is
\[
\D (\Theta_A(\delta_\chi))=\frac1{|A|}\sum_{a\in A}\chi(a^{-1})\,(a\otimes a).
\]
The right-hand side:
\[
(\Theta_A\otimes\Theta_A)\D (\delta_\chi)
=\sum_{\alpha\beta=\chi}\Theta_A(\delta_\alpha)\otimes\Theta_A(\delta_\beta)
=\frac1{|A|^2}\sum_{\alpha\beta=\chi}\sum_{a,b\in A}\alpha(a^{-1})\beta(b^{-1})\,(a\otimes b).
\]
Writing $\beta=\alpha^{-1}\chi$, the coefficient of a fixed $a\otimes b$ is
\[
\frac1{|A|^2}\sum_{\alpha\in A^\vee}\alpha(a^{-1})(\alpha^{-1}\chi)(b^{-1})
=\frac{\chi(b^{-1})}{|A|^2}\sum_{\alpha\in A^\vee}\alpha(a^{-1}b)
=\begin{cases}
\frac1{|A|}\chi(a^{-1}),& a=b,\\
0,& a\neq b,
\end{cases}
\]
which equals $\D (\Theta_A(\delta_\chi))$.

\medskip\noindent
{\it Counit and antipode.}
For the counit,
\[
\varepsilon(\Theta_A(\delta_\chi))=\frac1{|A|}\sum_{a\in A}\chi(a^{-1})=\delta_{\chi,1}=\varepsilon(\delta_\chi).
\]
For the antipode,
\[
S(\Theta_A(\delta_\chi))
=\frac1{|A|}\sum_{a\in A}\chi(a^{-1})\,a^{-1}
=\frac1{|A|}\sum_{a\in A}\chi(a)\,a
=\Theta_A(\delta_{\chi^{-1}})
=\Theta_A(S(\delta_\chi)).
\]



\medskip\noindent
{\it Inverse.}
We have show that $\Theta_A$ is a Hopf algebra map and hence it remains to verify the inverse given as in \Cref{eq:ThetaInv}.


Fix $a\in A$. Then by \eqref{eq:ThetaDef},
\[
\Theta_A(\Theta_A^{-1}(a))
=\sum_{\chi\in A^\vee}\chi(a)\,\Theta_A(\delta_\chi)
=\frac1{|A|}\sum_{\chi\in A^\vee}\ \sum_{x\in A}\chi(a)\chi(x^{-1})\,x
=\frac1{|A|}\sum_{x\in A}\Bigl(\sum_{\chi\in A^\vee}\chi(ax^{-1})\Bigr)x.
\]
By orthogonality \eqref{eq:orth}, $\sum_{\chi}\chi(ax^{-1})=|A|$ iff $x=a$ and $0$ otherwise. Hence the above equals $a$.
Conversely, for $\delta_\psi$ we compute using \eqref{eq:ThetaInv}:
\begin{align*}
\Theta_A^{-1}(\Theta_A(\delta_\psi))
& =\frac1{|A|}\sum_{x\in A}\psi(x^{-1})\,\Theta_A^{-1}(x)\\
& =\frac1{|A|}\sum_{x\in A}\psi(x^{-1})\sum_{\chi\in A^\vee}\chi(x)\delta_\chi \\
& =\sum_{\chi\in A^\vee}\Bigl(\frac1{|A|}\sum_{x\in A}(\chi\psi^{-1})(x)\Bigr)\delta_\chi
=\delta_\psi,
\end{align*}
again by \eqref{eq:orth}. Thus $\Theta_A^{-1}$ is indeed the inverse.
\end{proof}

\end{document}
